\documentclass[11pt,oneside]{book}

%%%%%%%%%%%%%%%%%%%%%%%%%%%%%%%%%%%%%%%%%%%%%%%%%%%%%%%%%%%%%%%%%%%%%%%%%%%%%%%%
% Packages
%%%%%%%%%%%%%%%%%%%%%%%%%%%%%%%%%%%%%%%%%%%%%%%%%%%%%%%%%%%%%%%%%%%%%%%%%%%%%%%%

\usepackage[T1]{fontenc}
\usepackage[utf8]{inputenc}
\usepackage[english]{babel}

\usepackage{lmodern}
\usepackage{microtype}

\usepackage{amsmath}
\usepackage{amssymb}
\usepackage{amsthm}
\usepackage{mathtools}
\usepackage{bm}

\usepackage{graphicx}
\usepackage{xcolor}
\usepackage{booktabs}
\usepackage{longtable}
\usepackage{array}
\usepackage{tabularx}
\usepackage{multirow}

\usepackage{enumitem}
\usepackage{csquotes}

\usepackage{tikz}
\usetikzlibrary{
	arrows.meta,
	positioning,
	shapes.geometric,
	fit,
	backgrounds,
	calc,
	matrix
}

\usepackage[
a4paper,
margin=2.8cm,
headheight=15pt
]{geometry}

\usepackage{fancyhdr}
\usepackage{hyperref}
\usepackage{bookmark}

%%%%%%%%%%%%%%%%%%%%%%%%%%%%%%%%%%%%%%%%%%%%%%%%%%%%%%%%%%%%%%%%%%%%%%%%%%%%%%%%
% Hyperlink configuration
%%%%%%%%%%%%%%%%%%%%%%%%%%%%%%%%%%%%%%%%%%%%%%%%%%%%%%%%%%%%%%%%%%%%%%%%%%%%%%%%

\hypersetup{
	colorlinks=true,
	linkcolor=blue!55!black,
	citecolor=green!40!black,
	urlcolor=blue!60!black,
	pdftitle={Entropy in Image Processing and Computer Vision},
	pdfauthor={Author},
	pdfsubject={Information theory, entropy, uncertainty, and computer vision},
	pdfkeywords={
		entropy,
		image processing,
		computer vision,
		rough sets,
		fuzzy sets,
		deep learning
	}
}

%%%%%%%%%%%%%%%%%%%%%%%%%%%%%%%%%%%%%%%%%%%%%%%%%%%%%%%%%%%%%%%%%%%%%%%%%%%%%%%%
% Page style
%%%%%%%%%%%%%%%%%%%%%%%%%%%%%%%%%%%%%%%%%%%%%%%%%%%%%%%%%%%%%%%%%%%%%%%%%%%%%%%%

\pagestyle{fancy}
\fancyhf{}
\fancyhead[L]{\nouppercase{\leftmark}}
\fancyhead[R]{\thepage}

\renewcommand{\headrulewidth}{0.4pt}
\renewcommand{\footrulewidth}{0pt}

%%%%%%%%%%%%%%%%%%%%%%%%%%%%%%%%%%%%%%%%%%%%%%%%%%%%%%%%%%%%%%%%%%%%%%%%%%%%%%%%
% Paragraph formatting
%%%%%%%%%%%%%%%%%%%%%%%%%%%%%%%%%%%%%%%%%%%%%%%%%%%%%%%%%%%%%%%%%%%%%%%%%%%%%%%%

\setlength{\parindent}{0pt}
\setlength{\parskip}{0.6em}

%%%%%%%%%%%%%%%%%%%%%%%%%%%%%%%%%%%%%%%%%%%%%%%%%%%%%%%%%%%%%%%%%%%%%%%%%%%%%%%%
% List formatting
%%%%%%%%%%%%%%%%%%%%%%%%%%%%%%%%%%%%%%%%%%%%%%%%%%%%%%%%%%%%%%%%%%%%%%%%%%%%%%%%

\setlist[itemize]{
	topsep=0.4em,
	itemsep=0.25em,
	parsep=0pt
}

\setlist[enumerate]{
	topsep=0.4em,
	itemsep=0.25em,
	parsep=0pt
}

%%%%%%%%%%%%%%%%%%%%%%%%%%%%%%%%%%%%%%%%%%%%%%%%%%%%%%%%%%%%%%%%%%%%%%%%%%%%%%%%
% Theorem environments
%%%%%%%%%%%%%%%%%%%%%%%%%%%%%%%%%%%%%%%%%%%%%%%%%%%%%%%%%%%%%%%%%%%%%%%%%%%%%%%%

\newtheorem{definition}{Definition}[chapter]
\newtheorem{theorem}{Theorem}[chapter]
\newtheorem{proposition}{Proposition}[chapter]
\newtheorem{lemma}{Lemma}[chapter]
\newtheorem{corollary}{Corollary}[chapter]

\theoremstyle{remark}
\newtheorem{remark}{Remark}[chapter]
\newtheorem{example}{Example}[chapter]

%%%%%%%%%%%%%%%%%%%%%%%%%%%%%%%%%%%%%%%%%%%%%%%%%%%%%%%%%%%%%%%%%%%%%%%%%%%%%%%%
% Custom commands
%%%%%%%%%%%%%%%%%%%%%%%%%%%%%%%%%%%%%%%%%%%%%%%%%%%%%%%%%%%%%%%%%%%%%%%%%%%%%%%%

\newcommand{\Entropy}{\mathcal{H}}
\newcommand{\Prob}{\mathbb{P}}
\newcommand{\Expectation}{\mathbb{E}}
\newcommand{\Lower}{\underline{\operatorname{apr}}}
\newcommand{\Upper}{\overline{\operatorname{apr}}}
\newcommand{\Boundary}{\operatorname{BND}}
\newcommand{\Positive}{\operatorname{POS}}
\newcommand{\Negative}{\operatorname{NEG}}
\newcommand{\KL}{D_{\mathrm{KL}}}
\newcommand{\JS}{D_{\mathrm{JS}}}
\newcommand{\R}{\mathbb{R}}

%%%%%%%%%%%%%%%%%%%%%%%%%%%%%%%%%%%%%%%%%%%%%%%%%%%%%%%%%%%%%%%%%%%%%%%%%%%%%%%%
% Custom environments
%%%%%%%%%%%%%%%%%%%%%%%%%%%%%%%%%%%%%%%%%%%%%%%%%%%%%%%%%%%%%%%%%%%%%%%%%%%%%%%%

\newenvironment{chapteroverview}
{
	\begin{quote}
		\small
		\textbf{Chapter overview.}
	}
	{
	\end{quote}
}

\newenvironment{researchquestions}
{
	\begin{quote}
		\small
		\textbf{Research questions.}
		\begin{enumerate}
		}
		{
		\end{enumerate}
	\end{quote}
}

\newenvironment{datasetlist}
{
	\begin{itemize}[label=\(\triangleright\)]
	}
	{
	\end{itemize}
}

\newenvironment{applicationlist}
{
	\begin{itemize}[label=\(\blacktriangleright\)]
	}
	{
	\end{itemize}
}

%%%%%%%%%%%%%%%%%%%%%%%%%%%%%%%%%%%%%%%%%%%%%%%%%%%%%%%%%%%%%%%%%%%%%%%%%%%%%%%%
% Document metadata
%%%%%%%%%%%%%%%%%%%%%%%%%%%%%%%%%%%%%%%%%%%%%%%%%%%%%%%%%%%%%%%%%%%%%%%%%%%%%%%%

\title{
	\Huge
	Entropy in Image Processing\\
	and Computer Vision
	\\[1em]
	\Large
	From Classical Information Theory to\\
	Rough, Fuzzy, Graph-Based, and Deep Entropy
}

\author{Author Name}

\date{\today}

%%%%%%%%%%%%%%%%%%%%%%%%%%%%%%%%%%%%%%%%%%%%%%%%%%%%%%%%%%%%%%%%%%%%%%%%%%%%%%%%
% Document
%%%%%%%%%%%%%%%%%%%%%%%%%%%%%%%%%%%%%%%%%%%%%%%%%%%%%%%%%%%%%%%%%%%%%%%%%%%%%%%%

\begin{document}
	
	%%%%%%%%%%%%%%%%%%%%%%%%%%%%%%%%%%%%%%%%%%%%%%%%%%%%%%%%%%%%%%%%%%%%%%%%%%%%%%%%
	% Front matter
	%%%%%%%%%%%%%%%%%%%%%%%%%%%%%%%%%%%%%%%%%%%%%%%%%%%%%%%%%%%%%%%%%%%%%%%%%%%%%%%%
	
	\frontmatter
	
	\maketitle
	
	\chapter*{Preface}
	\addcontentsline{toc}{chapter}{Preface}
	
	Entropy has long served as a fundamental measure of information, uncertainty,
	complexity, and disorder. In classical image processing, entropy is commonly
	computed from image-intensity histograms or local pixel neighborhoods. Modern
	computer vision, however, operates on increasingly abstract representations,
	including superpixels, region graphs, convolutional feature maps, transformer
	tokens, and high-dimensional latent embeddings.
	
	This book develops a unified account of entropy across these representations.
	It connects classical information theory with probabilistic uncertainty,
	fuzzy sets, rough sets, fuzzy-rough systems, evidence theory, graph methods,
	deep learning, Vision Transformers, and foundation models.
	
	The principal thesis is that entropy should be regarded not as a single
	formula applied exclusively to pixel histograms, but as a general functional
	defined over a representation and its associated model of uncertainty.
	
	\chapter*{Notation}
	\addcontentsline{toc}{chapter}{Notation}
	
	\begin{longtable}{p{0.22\textwidth}p{0.68\textwidth}}
		\toprule
		Symbol & Meaning \\
		\midrule
		\endfirsthead
		
		\toprule
		Symbol & Meaning \\
		\midrule
		\endhead
		
		\(X,Y,Z\) &
		Random variables or learned representations.
		\\
		
		\(p(x)\) &
		Probability mass or density function.
		\\
		
		\(\Entropy(X)\) &
		Entropy of \(X\).
		\\
		
		\(\Lower{R}(A)\) &
		Lower rough approximation of \(A\) under relation \(R\).
		\\
		
		\(\Upper{R}(A)\) &
		Upper rough approximation of \(A\) under relation \(R\).
		\\
		
		\(\Boundary_R(A)\) &
		Rough boundary region of \(A\).
		\\
		
		\(\mu_A(x)\) &
		Fuzzy membership of \(x\) in \(A\).
		\\
		
		\(G=(V,E)\) &
		Graph with node set \(V\) and edge set \(E\).
		\\
		
		\(\bm{z}_i\) &
		Deep feature or latent embedding of sample \(i\).
		\\
		
		\bottomrule
	\end{longtable}
	
	\tableofcontents
	\listoffigures
	\listoftables
	
	%%%%%%%%%%%%%%%%%%%%%%%%%%%%%%%%%%%%%%%%%%%%%%%%%%%%%%%%%%%%%%%%%%%%%%%%%%%%%%%%
	% Main matter
	%%%%%%%%%%%%%%%%%%%%%%%%%%%%%%%%%%%%%%%%%%%%%%%%%%%%%%%%%%%%%%%%%%%%%%%%%%%%%%%%
	
	\mainmatter
	
	%%%%%%%%%%%%%%%%%%%%%%%%%%%%%%%%%%%%%%%%%%%%%%%%%%%%%%%%%%%%%%%%%%%%%%%%%%%%%%%%
	% PART I
	%%%%%%%%%%%%%%%%%%%%%%%%%%%%%%%%%%%%%%%%%%%%%%%%%%%%%%%%%%%%%%%%%%%%%%%%%%%%%%%%
	
	\part{Foundations}
	
	%%%%%%%%%%%%%%%%%%%%%%%%%%%%%%%%%%%%%%%%%%%%%%%%%%%%%%%%%%%%%%%%%%%%%%%%%%%%%%%%
	% Chapter 1
	%%%%%%%%%%%%%%%%%%%%%%%%%%%%%%%%%%%%%%%%%%%%%%%%%%%%%%%%%%%%%%%%%%%%%%%%%%%%%%%%
	
	\chapter{Introduction}
	\label{chap:introduction}
	
	\begin{chapteroverview}
		This chapter introduces information, uncertainty, and entropy in image
		processing. It presents the historical development of entropy methods,
		defines the scope of the book, and explains its organization.
	\end{chapteroverview}
	
	\section{Information and Uncertainty}
	
	Information quantifies how much is learned when an observation is made.
	Uncertainty describes the lack of complete knowledge before that observation.
	Entropy connects these ideas by assigning a numerical value to the uncertainty
	of a probability distribution or, more generally, to the ambiguity of a
	mathematical representation.
	
	\section{Why Entropy in Image Processing?}
	
	Entropy has been used in image processing for:
	
	\begin{itemize}
		\item image thresholding;
		\item texture characterization;
		\item image registration;
		\item compression;
		\item segmentation;
		\item feature selection;
		\item uncertainty estimation;
		\item active learning;
		\item representation analysis.
	\end{itemize}
	
	\section{Historical Development}
	
	The development of entropy methods may be viewed as a transition through
	successively richer representations:



	\begin{center}
		\begin{tikzpicture}[
			node distance=8mm,
			box/.style={
				draw,
				rounded corners,
				minimum width=7cm,
				minimum height=8mm,
				align=center,
				fill=blue!5
			},
			arrow/.style={
				-{Latex[length=2.5mm]},
				thick
			}
			]
			
			\node[box] (pixels) {Pixel intensities};
			\node[box, below=of pixels] (local) {Local neighborhoods and textures};
			\node[box, below=of local] (regions) {Regions and superpixels};
			\node[box, below=of regions] (graphs) {Graphs and relational structures};
			\node[box, below=of graphs] (features) {Deep feature maps and embeddings};
			\node[box, below=of features] (foundation)
			{Foundation-model representations};
			
			\draw[arrow] (pixels) -- (local);
			\draw[arrow] (local) -- (regions);
			\draw[arrow] (regions) -- (graphs);
			\draw[arrow] (graphs) -- (features);
			\draw[arrow] (features) -- (foundation);
			
		\end{tikzpicture}
	\end{center}
	
	\section{Scope of the Survey}
	
	This book covers five interacting dimensions:
	
	\begin{enumerate}
		\item image representations;
		\item mathematical uncertainty models;
		\item entropy measures;
		\item learning architectures;
		\item application domains.
	\end{enumerate}
	
	\section{Organization of the Book}
	
	The book is divided into six parts:
	
	\begin{description}
		\item[Part I:] Foundations of information theory and generalized entropy.
		\item[Part II:] Entropy on image representations.
		\item[Part III:] Probabilistic, fuzzy, rough, and evidence-based uncertainty.
		\item[Part IV:] Entropy in modern deep learning.
		\item[Part V:] Application domains and benchmark datasets.
		\item[Part VI:] A unified framework and future research directions.
	\end{description}
	
	%%%%%%%%%%%%%%%%%%%%%%%%%%%%%%%%%%%%%%%%%%%%%%%%%%%%%%%%%%%%%%%%%%%%%%%%%%%%%%%%
	% Chapter 2
	%%%%%%%%%%%%%%%%%%%%%%%%%%%%%%%%%%%%%%%%%%%%%%%%%%%%%%%%%%%%%%%%%%%%%%%%%%%%%%%%
	
	\chapter{Information Theory}
	\label{chap:information-theory}
	
	\begin{chapteroverview}
		This chapter develops the mathematical foundations required throughout the
		book, including Shannon entropy, mutual information, divergences, and the
		information bottleneck.
	\end{chapteroverview}
	
	\section{Shannon Entropy}
	
	For a discrete random variable \(X\) with outcomes
	\(\{x_1,\ldots,x_n\}\) and probabilities \(p_i=\Prob(X=x_i)\),
	Shannon entropy is
	
	\[
	\Entropy(X)
	=
	-\sum_{i=1}^{n} p_i \log p_i.
	\]
	
	\section{Joint Entropy}
	
	For random variables \(X\) and \(Y\),
	
	\[
	\Entropy(X,Y)
	=
	-\sum_{x}\sum_{y}
	p(x,y)\log p(x,y).
	\]
	
	\section{Conditional Entropy}
	
	Conditional entropy is defined by
	
	\[
	\Entropy(Y\mid X)
	=
	\Entropy(X,Y)-\Entropy(X).
	\]
	
	\section{Mutual Information}
	
	Mutual information measures the statistical dependence between \(X\) and \(Y\):
	
	\[
	I(X;Y)
	=
	\Entropy(X)+\Entropy(Y)-\Entropy(X,Y).
	\]
	
	Equivalently,
	
	\[
	I(X;Y)
	=
	\sum_{x,y}
	p(x,y)
	\log
	\frac{p(x,y)}{p(x)p(y)}.
	\]
	
	\section{Relative Entropy}
	
	The Kullback--Leibler divergence between distributions \(P\) and \(Q\) is
	
	\[
	\KL(P\Vert Q)
	=
	\sum_x p(x)
	\log
	\frac{p(x)}{q(x)}.
	\]
	
	\section{Cross Entropy}
	
	The cross entropy between \(P\) and \(Q\) is
	
	\[
	\Entropy(P,Q)
	=
	-\sum_x p(x)\log q(x).
	\]
	
	\section{Jensen--Shannon Divergence}
	
	The Jensen--Shannon divergence is
	
	\[
	\JS(P,Q)
	=
	\frac{1}{2}
	\KL(P\Vert M)
	+
	\frac{1}{2}
	\KL(Q\Vert M),
	\]
	
	where
	
	\[
	M=\frac{1}{2}(P+Q).
	\]
	
	\section{Information Bottleneck}
	
	The information bottleneck objective seeks a representation \(Z\) that
	preserves information relevant to target \(Y\) while compressing input \(X\):
	
	\[
	\mathcal{L}_{\mathrm{IB}}
	=
	I(X;Z)-\beta I(Z;Y).
	\]
	
	%%%%%%%%%%%%%%%%%%%%%%%%%%%%%%%%%%%%%%%%%%%%%%%%%%%%%%%%%%%%%%%%%%%%%%%%%%%%%%%%
	% Chapter 3
	%%%%%%%%%%%%%%%%%%%%%%%%%%%%%%%%%%%%%%%%%%%%%%%%%%%%%%%%%%%%%%%%%%%%%%%%%%%%%%%%
	
	\chapter{Generalized Entropies}
	\label{chap:generalized-entropies}
	
	\begin{chapteroverview}
		This chapter introduces entropy families that extend Shannon entropy or measure
		regularity, predictability, and complexity in ordered and multiscale data.
	\end{chapteroverview}
	
	\section{Rényi Entropy}
	
	For order \(\alpha>0\), \(\alpha\neq1\),
	
	\[
	\Entropy_{\alpha}^{\mathrm{R}}(X)
	=
	\frac{1}{1-\alpha}
	\log
	\left(
	\sum_i p_i^\alpha
	\right).
	\]
	
	\section{Tsallis Entropy}
	
	For \(q>0\), \(q\neq1\),
	
	\[
	\Entropy_q^{\mathrm{T}}(X)
	=
	\frac{1}{q-1}
	\left(
	1-\sum_i p_i^q
	\right).
	\]
	
	\section{Havrda--Charvát Entropy}
	
	The Havrda--Charvát entropy belongs to the family of generalized
	nonadditive entropy measures. Its relationship to the Tsallis family should
	be discussed together with differences in normalization conventions.
	
	\section{Sharma--Mittal Entropy}
	
	The Sharma--Mittal family combines parameters associated with Rényi and
	Tsallis entropies and provides a broad two-parameter entropy model.
	
	\section{Approximate Entropy}
	
	Approximate entropy measures the regularity and predictability of ordered data.
	It is particularly useful for texture sequences, biomedical signals, and
	spatially ordered image descriptors.
	
	\section{Sample Entropy}
	
	Sample entropy modifies approximate entropy by excluding self-matches and
	usually provides improved statistical behavior for finite sequences.
	
	\section{Permutation Entropy}
	
	Permutation entropy measures the diversity of ordinal patterns in a sequence
	or local image neighborhood.
	
	\section{Multiscale Entropy}
	
	Multiscale entropy evaluates complexity across several spatial or temporal
	scales rather than at a single resolution.
	
	%%%%%%%%%%%%%%%%%%%%%%%%%%%%%%%%%%%%%%%%%%%%%%%%%%%%%%%%%%%%%%%%%%%%%%%%%%%%%%%%
	% PART II
	%%%%%%%%%%%%%%%%%%%%%%%%%%%%%%%%%%%%%%%%%%%%%%%%%%%%%%%%%%%%%%%%%%%%%%%%%%%%%%%%
	
	\part{Entropy on Image Representations}
	
	%%%%%%%%%%%%%%%%%%%%%%%%%%%%%%%%%%%%%%%%%%%%%%%%%%%%%%%%%%%%%%%%%%%%%%%%%%%%%%%%
	% Chapter 4
	%%%%%%%%%%%%%%%%%%%%%%%%%%%%%%%%%%%%%%%%%%%%%%%%%%%%%%%%%%%%%%%%%%%%%%%%%%%%%%%%
	
	\chapter{Pixel-Based Entropy}
	\label{chap:pixel-entropy}
	
	\begin{chapteroverview}
		Pixel-based entropy treats image intensities or local intensity distributions
		as the primary representation.
	\end{chapteroverview}
	
	\section{Histogram Entropy}
	
	Let an image have normalized intensity histogram
	\(\{p_0,\ldots,p_{L-1}\}\). Its global histogram entropy is
	
	\[
	\Entropy_{\mathrm{hist}}
	=
	-\sum_{i=0}^{L-1} p_i\log p_i.
	\]
	
	\section{Local Entropy}
	
	For a neighborhood \(\mathcal{N}(x)\) around pixel \(x\),
	
	\[
	\Entropy_{\mathrm{local}}(x)
	=
	-\sum_{k=0}^{L-1}
	p_k\bigl(\mathcal{N}(x)\bigr)
	\log
	p_k\bigl(\mathcal{N}(x)\bigr).
	\]
	
	\section{Multiscale Entropy}
	
	Local entropy can be evaluated using several neighborhood sizes or image
	pyramid levels.
	
	\section{Adaptive Entropy}
	
	Adaptive methods choose neighborhood size, histogram binning, or entropy
	parameters according to local image structure.
	
	\section{Thresholding}
	
	Entropy-based thresholding chooses threshold \(t\) by optimizing an objective
	defined over foreground and background distributions.
	
	\section*{Representative Datasets}
	\addcontentsline{toc}{section}{Representative Datasets}
	
	\begin{datasetlist}
		\item Kodak image dataset;
		\item Berkeley Segmentation Dataset and Benchmark, BSDS500.
	\end{datasetlist}
	
	%%%%%%%%%%%%%%%%%%%%%%%%%%%%%%%%%%%%%%%%%%%%%%%%%%%%%%%%%%%%%%%%%%%%%%%%%%%%%%%%
	% Chapter 5
	%%%%%%%%%%%%%%%%%%%%%%%%%%%%%%%%%%%%%%%%%%%%%%%%%%%%%%%%%%%%%%%%%%%%%%%%%%%%%%%%
	
	\chapter{Texture Entropy}
	\label{chap:texture-entropy}
	
	\begin{chapteroverview}
		Texture entropy quantifies local spatial organization, repetition,
		irregularity, and multiresolution structure.
	\end{chapteroverview}
	
	\section{GLCM Entropy}
	
	For a normalized gray-level co-occurrence matrix \(P(i,j)\),
	
	\[
	\Entropy_{\mathrm{GLCM}}
	=
	-\sum_i\sum_j
	P(i,j)\log P(i,j).
	\]
	
	\section{Wavelet Entropy}
	
	Wavelet entropy measures the distribution of energy among wavelet
	coefficients or frequency subbands.
	
	\section{Fractal Entropy}
	
	Fractal entropy combines scale-dependent structural complexity with
	fractal representations.
	
	\section{Local Binary Patterns}
	
	Entropy may be computed over local binary pattern histograms to measure the
	diversity of local texture structures.
	
	\section{Multiresolution Entropy}
	
	Multiresolution methods analyze entropy across scales, frequency bands, or
	wavelet decomposition levels.
	
	\section*{Representative Datasets}
	\addcontentsline{toc}{section}{Representative Datasets}
	
	\begin{datasetlist}
		\item Brodatz texture dataset;
		\item Describable Textures Dataset;
		\item KTH-TIPS texture dataset.
	\end{datasetlist}
	
	%%%%%%%%%%%%%%%%%%%%%%%%%%%%%%%%%%%%%%%%%%%%%%%%%%%%%%%%%%%%%%%%%%%%%%%%%%%%%%%%
	% Chapter 6
	%%%%%%%%%%%%%%%%%%%%%%%%%%%%%%%%%%%%%%%%%%%%%%%%%%%%%%%%%%%%%%%%%%%%%%%%%%%%%%%%
	
	\chapter{Region Entropy}
	\label{chap:region-entropy}
	
	\begin{chapteroverview}
		Region entropy moves beyond individual pixels and measures uncertainty,
		heterogeneity, or structural complexity within image regions.
	\end{chapteroverview}
	
	\section{Superpixels}
	
	Superpixels group nearby pixels with similar appearance into perceptually
	meaningful atomic regions.
	
	\section{Region Descriptors}
	
	A region may be represented by:
	
	\begin{itemize}
		\item mean color;
		\item texture statistics;
		\item shape descriptors;
		\item spatial coordinates;
		\item deep feature embeddings;
		\item class-probability distributions.
	\end{itemize}
	
	\section{Region Uncertainty}
	
	For region \(R\), one possible uncertainty measure is
	
	\[
	\Entropy(R)
	=
	-\sum_{c=1}^{C}
	p(c\mid R)\log p(c\mid R).
	\]
	
	\section{Region Merging}
	
	Entropy may guide region merging by measuring whether merging two adjacent
	regions reduces internal heterogeneity or increases semantic consistency.
	
	\section{Region Entropy}
	
	Region entropy may incorporate class uncertainty, internal feature variation,
	boundary uncertainty, and relations with adjacent regions.
	
	\section*{Representative Datasets}
	\addcontentsline{toc}{section}{Representative Datasets}
	
	\begin{datasetlist}
		\item Cityscapes;
		\item ADE20K.
	\end{datasetlist}
	
	%%%%%%%%%%%%%%%%%%%%%%%%%%%%%%%%%%%%%%%%%%%%%%%%%%%%%%%%%%%%%%%%%%%%%%%%%%%%%%%%
	% Chapter 7
	%%%%%%%%%%%%%%%%%%%%%%%%%%%%%%%%%%%%%%%%%%%%%%%%%%%%%%%%%%%%%%%%%%%%%%%%%%%%%%%%
	
	\chapter{Graph Entropy}
	\label{chap:graph-entropy}
	
	\begin{chapteroverview}
		This chapter treats images as graphs whose nodes represent pixels,
		superpixels, regions, objects, or learned tokens.
	\end{chapteroverview}
	
	\section{Image Graphs}
	
	An image graph is written as
	
	\[
	G=(V,E),
	\]
	
	where \(V\) is the set of nodes and \(E\) represents adjacency, similarity,
	or semantic relationships.
	
	\section{Region Adjacency Graphs}
	
	A region adjacency graph connects spatially adjacent superpixels or segmented
	regions.
	
	\section{Graph Laplacian}
	
	Given adjacency matrix \(A\) and degree matrix \(D\), the graph Laplacian is
	
	\[
	L=D-A.
	\]
	
	The normalized graph Laplacian may be written as
	
	\[
	L_{\mathrm{norm}}
	=
	I-D^{-1/2}AD^{-1/2}.
	\]
	
	\section{Structural Entropy}
	
	Structural entropy measures uncertainty associated with the organization,
	partitioning, or coding structure of a graph.
	
	\section{Spectral Entropy}
	
	Let \(\lambda_1,\ldots,\lambda_n\) denote normalized spectral quantities such
	that \(\sum_i\lambda_i=1\). Spectral entropy may be defined as
	
	\[
	\Entropy_{\mathrm{spectral}}
	=
	-\sum_{i=1}^{n}
	\lambda_i\log\lambda_i.
	\]
	
	\section{Random-Walk Entropy}
	
	Random-walk entropy measures uncertainty in graph transitions or trajectories.
	
	\section*{Representative Datasets}
	\addcontentsline{toc}{section}{Representative Datasets}
	
	\begin{datasetlist}
		\item Microsoft COCO;
		\item PASCAL VOC.
	\end{datasetlist}
	
	%%%%%%%%%%%%%%%%%%%%%%%%%%%%%%%%%%%%%%%%%%%%%%%%%%%%%%%%%%%%%%%%%%%%%%%%%%%%%%%%
	% Chapter 8
	%%%%%%%%%%%%%%%%%%%%%%%%%%%%%%%%%%%%%%%%%%%%%%%%%%%%%%%%%%%%%%%%%%%%%%%%%%%%%%%%
	
	\chapter{Feature Entropy}
	\label{chap:feature-entropy}
	
	\begin{chapteroverview}
		Feature entropy measures uncertainty or diversity in representations learned
		by neural networks.
	\end{chapteroverview}
	
	\section{CNN Feature Maps}
	
	For an activation tensor
	
	\[
	F\in\R^{H\times W\times D},
	\]
	
	each spatial location may be represented by a \(D\)-dimensional vector.
	
	\section{Deep Embeddings}
	
	A neural encoder maps input image \(x_i\) to embedding
	
	\[
	\bm{z}_i=f_{\theta}(x_i).
	\]
	
	\section{Latent Representations}
	
	Latent entropy may describe dispersion, compression, class structure,
	or information content in the learned space.
	
	\section{Representation Entropy}
	
	Representation entropy may be estimated using:
	
	\begin{itemize}
		\item histograms;
		\item kernel density estimation;
		\item nearest-neighbor estimators;
		\item covariance spectra;
		\item graph-based measures;
		\item matrix-based Rényi entropy.
	\end{itemize}
	
	\section{Attention Entropy}
	
	For attention distribution \(a_{ij}\),
	
	\[
	\Entropy_{\mathrm{attn}}(i)
	=
	-\sum_j a_{ij}\log a_{ij}.
	\]
	
	\section*{Representative Models and Datasets}
	\addcontentsline{toc}{section}{Representative Models and Datasets}
	
	\begin{datasetlist}
		\item ImageNet;
		\item Segment Anything representations;
		\item DINO representations.
	\end{datasetlist}
	
	%%%%%%%%%%%%%%%%%%%%%%%%%%%%%%%%%%%%%%%%%%%%%%%%%%%%%%%%%%%%%%%%%%%%%%%%%%%%%%%%
	% PART III
	%%%%%%%%%%%%%%%%%%%%%%%%%%%%%%%%%%%%%%%%%%%%%%%%%%%%%%%%%%%%%%%%%%%%%%%%%%%%%%%%
	
	\part{Uncertainty Models}
	
	%%%%%%%%%%%%%%%%%%%%%%%%%%%%%%%%%%%%%%%%%%%%%%%%%%%%%%%%%%%%%%%%%%%%%%%%%%%%%%%%
	% Chapter 9
	%%%%%%%%%%%%%%%%%%%%%%%%%%%%%%%%%%%%%%%%%%%%%%%%%%%%%%%%%%%%%%%%%%%%%%%%%%%%%%%%
	
	\chapter{Probabilistic Entropy}
	\label{chap:probabilistic-entropy}
	
	\begin{chapteroverview}
		Probabilistic entropy describes uncertainty using probability distributions
		over labels, model parameters, or predictions.
	\end{chapteroverview}
	
	\section{Bayesian Inference}
	
	Bayesian inference updates uncertainty about parameters \(\theta\) using data
	\(\mathcal{D}\):
	
	\[
	p(\theta\mid\mathcal{D})
	=
	\frac{
		p(\mathcal{D}\mid\theta)p(\theta)
	}{
		p(\mathcal{D})
	}.
	\]
	
	\section{Predictive Entropy}
	
	For class probabilities \(p(y=c\mid x)\),
	
	\[
	\Entropy_{\mathrm{pred}}(x)
	=
	-\sum_{c=1}^{C}
	p(y=c\mid x)
	\log p(y=c\mid x).
	\]
	
	\section{Ensemble Entropy}
	
	An ensemble combines predictions from several models:
	
	\[
	\bar{p}(y\mid x)
	=
	\frac{1}{M}
	\sum_{m=1}^{M}
	p_m(y\mid x).
	\]
	
	Entropy is then computed from \(\bar{p}(y\mid x)\).
	
	\section{Monte Carlo Dropout}
	
	Repeated stochastic forward passes provide an approximate predictive
	distribution.
	
	\section{Calibration}
	
	A calibrated model should assign confidence values that correspond to empirical
	correctness frequencies.
	
	%%%%%%%%%%%%%%%%%%%%%%%%%%%%%%%%%%%%%%%%%%%%%%%%%%%%%%%%%%%%%%%%%%%%%%%%%%%%%%%%
	% Chapter 10
	%%%%%%%%%%%%%%%%%%%%%%%%%%%%%%%%%%%%%%%%%%%%%%%%%%%%%%%%%%%%%%%%%%%%%%%%%%%%%%%%
	
	\chapter{Fuzzy Entropy}
	\label{chap:fuzzy-entropy}
	
	\begin{chapteroverview}
		Fuzzy entropy quantifies vagueness represented by gradual memberships rather
		than crisp class assignments.
	\end{chapteroverview}
	
	\section{Fuzzy Sets}
	
	A fuzzy set \(A\) on universe \(U\) is characterized by membership function
	
	\[
	\mu_A:U\rightarrow[0,1].
	\]
	
	\section{Membership Functions}
	
	Common membership functions include triangular, trapezoidal, Gaussian,
	sigmoidal, and data-driven forms.
	
	\section{Fuzzy Entropy}
	
	A generic fuzzy entropy may be defined as
	
	\[
	\Entropy_F(A)
	=
	\frac{1}{|U|}
	\sum_{x\in U}
	h\bigl(\mu_A(x)\bigr),
	\]
	
	where \(h\) is maximal near \(\mu_A(x)=1/2\) and minimal near \(0\) or \(1\).
	
	\section{Type-2 Fuzzy Entropy}
	
	Type-2 fuzzy systems assign fuzzy uncertainty to the membership degree itself.
	
	\section{Intuitionistic Fuzzy Entropy}
	
	An intuitionistic fuzzy set specifies membership, non-membership, and
	hesitation.
	
	\section*{Applications}
	\addcontentsline{toc}{section}{Applications}
	
	\begin{applicationlist}
		\item medical image analysis;
		\item image segmentation;
		\item feature selection.
	\end{applicationlist}
	
	%%%%%%%%%%%%%%%%%%%%%%%%%%%%%%%%%%%%%%%%%%%%%%%%%%%%%%%%%%%%%%%%%%%%%%%%%%%%%%%%
	% Chapter 11
	%%%%%%%%%%%%%%%%%%%%%%%%%%%%%%%%%%%%%%%%%%%%%%%%%%%%%%%%%%%%%%%%%%%%%%%%%%%%%%%%
	
	\chapter{Rough Entropy}
	\label{chap:rough-entropy}
	
	\begin{chapteroverview}
		Rough entropy measures uncertainty caused by indiscernibility and imperfect
		approximation.
	\end{chapteroverview}
	
	\section{Rough Sets}
	
	Let \(U\) be a finite universe and let \(R\) be an equivalence relation or
	neighborhood relation on \(U\).
	
	\section{Lower Approximation}
	
	For \(A\subseteq U\), the lower approximation contains objects that certainly
	belong to \(A\):
	
	\[
	\Lower{R}(A)
	=
	\left\{
	x\in U:
	[x]_R\subseteq A
	\right\}.
	\]
	
	\section{Upper Approximation}
	
	The upper approximation contains objects that possibly belong to \(A\):
	
	\[
	\Upper{R}(A)
	=
	\left\{
	x\in U:
	[x]_R\cap A\neq\varnothing
	\right\}.
	\]
	
	\section{Boundary Region}
	
	The rough boundary is
	
	\[
	\Boundary_R(A)
	=
	\Upper{R}(A)
	\setminus
	\Lower{R}(A).
	\]
	
	\section{Rough Information Entropy}
	
	A rough entropy measure may quantify the size, diversity, or granularity of
	equivalence classes and approximation regions.
	
	\section{Neighborhood Rough Entropy}
	
	Neighborhood rough sets replace exact equivalence classes with metric or
	similarity neighborhoods.
	
	\section{Covering Rough Entropy}
	
	Covering rough sets permit overlapping granules and provide a more flexible
	representation than strict partitions.
	
	\section*{Applications}
	\addcontentsline{toc}{section}{Applications}
	
	\begin{applicationlist}
		\item hyperspectral imaging;
		\item remote sensing;
		\item medical imaging;
		\item feature selection;
		\item segmentation uncertainty.
	\end{applicationlist}
	
	%%%%%%%%%%%%%%%%%%%%%%%%%%%%%%%%%%%%%%%%%%%%%%%%%%%%%%%%%%%%%%%%%%%%%%%%%%%%%%%%
	% Chapter 12
	%%%%%%%%%%%%%%%%%%%%%%%%%%%%%%%%%%%%%%%%%%%%%%%%%%%%%%%%%%%%%%%%%%%%%%%%%%%%%%%%
	
	\chapter{Fuzzy-Rough Entropy}
	\label{chap:fuzzy-rough-entropy}
	
	\begin{chapteroverview}
		Fuzzy-rough entropy combines gradual membership with approximation-based
		uncertainty.
	\end{chapteroverview}
	
	\section{Hybrid Uncertainty}
	
	Fuzzy uncertainty represents gradual membership, whereas rough uncertainty
	represents indiscernibility between objects.
	
	\section{Fuzzy Neighborhoods}
	
	A fuzzy relation
	
	\[
	R:U\times U\rightarrow[0,1]
	\]
	
	assigns a similarity degree to every pair of objects.
	
	\section{Fuzzy-Rough Approximations}
	
	Fuzzy-rough lower and upper approximations combine fuzzy implication,
	aggregation, and neighborhood similarity.
	
	\section{Hybrid Entropy Measures}
	
	A fuzzy-rough entropy may combine:
	
	\begin{itemize}
		\item uncertainty of fuzzy memberships;
		\item uncertainty of lower and upper approximations;
		\item fuzzy boundary size;
		\item neighborhood inconsistency.
	\end{itemize}
	
	\section{Existing Algorithms}
	
	This section should classify algorithms according to:
	
	\begin{itemize}
		\item feature selection;
		\item clustering;
		\item classification;
		\item thresholding;
		\item segmentation;
		\item deep representation learning.
	\end{itemize}
	
	%%%%%%%%%%%%%%%%%%%%%%%%%%%%%%%%%%%%%%%%%%%%%%%%%%%%%%%%%%%%%%%%%%%%%%%%%%%%%%%%
	% Chapter 13
	%%%%%%%%%%%%%%%%%%%%%%%%%%%%%%%%%%%%%%%%%%%%%%%%%%%%%%%%%%%%%%%%%%%%%%%%%%%%%%%%
	
	\chapter{Evidence-Theoretic Entropy}
	\label{chap:evidence-entropy}
	
	\begin{chapteroverview}
		Evidence-theoretic entropy quantifies uncertainty represented by belief masses
		assigned to subsets of hypotheses.
	\end{chapteroverview}
	
	\section{Dempster--Shafer Theory}
	
	Let \(\Omega\) be a frame of discernment. A basic probability assignment is
	
	\[
	m:2^\Omega\rightarrow[0,1],
	\]
	
	such that
	
	\[
	m(\varnothing)=0,
	\qquad
	\sum_{A\subseteq\Omega}m(A)=1.
	\]
	
	\section{Belief Entropy}
	
	Belief entropy measures uncertainty in basic probability assignments.
	
	\section{Deng Entropy}
	
	Deng entropy is defined by
	
	\[
	\Entropy_{\mathrm{Deng}}(m)
	=
	-\sum_{A\subseteq\Omega}
	m(A)
	\log
	\left(
	\frac{m(A)}{2^{|A|}-1}
	\right).
	\]
	
	\section{Conflict Measures}
	
	Conflict measures quantify disagreement between evidence sources.
	
	\section{Sensor Fusion}
	
	Evidence theory is useful when combining predictions from cameras, LiDAR,
	radar, thermal sensors, and independent neural networks.
	
	%%%%%%%%%%%%%%%%%%%%%%%%%%%%%%%%%%%%%%%%%%%%%%%%%%%%%%%%%%%%%%%%%%%%%%%%%%%%%%%%
	% PART IV
	%%%%%%%%%%%%%%%%%%%%%%%%%%%%%%%%%%%%%%%%%%%%%%%%%%%%%%%%%%%%%%%%%%%%%%%%%%%%%%%%
	
	\part{Deep Learning}
	
	%%%%%%%%%%%%%%%%%%%%%%%%%%%%%%%%%%%%%%%%%%%%%%%%%%%%%%%%%%%%%%%%%%%%%%%%%%%%%%%%
	% Chapter 14
	%%%%%%%%%%%%%%%%%%%%%%%%%%%%%%%%%%%%%%%%%%%%%%%%%%%%%%%%%%%%%%%%%%%%%%%%%%%%%%%%
	
	\chapter{Entropy in Convolutional Neural Networks}
	\label{chap:cnn-entropy}
	
	\begin{chapteroverview}
		This chapter examines entropy in convolutional feature maps, activation
		patterns, deep representations, and self-supervised learning.
	\end{chapteroverview}
	
	\section{Feature Entropy}
	
	Feature entropy may be computed per channel, spatial location, layer, sample,
	or class.
	
	\section{Activation Entropy}
	
	Activation entropy measures the diversity and utilization of neural
	activations.
	
	\section{Representation Learning}
	
	Entropy can regulate the geometry and information content of learned
	embeddings.
	
	\section{Self-Supervised Learning}
	
	Entropy-based objectives may prevent representation collapse while preserving
	semantic invariance.
	
	%%%%%%%%%%%%%%%%%%%%%%%%%%%%%%%%%%%%%%%%%%%%%%%%%%%%%%%%%%%%%%%%%%%%%%%%%%%%%%%%
	% Chapter 15
	%%%%%%%%%%%%%%%%%%%%%%%%%%%%%%%%%%%%%%%%%%%%%%%%%%%%%%%%%%%%%%%%%%%%%%%%%%%%%%%%
	
	\chapter{Entropy in Vision Transformers}
	\label{chap:vit-entropy}
	
	\begin{chapteroverview}
		This chapter studies entropy in patch tokens, attention matrices, transformer
		representations, and semantic predictions.
	\end{chapteroverview}
	
	\section{Token Entropy}
	
	Token entropy measures diversity or uncertainty among patch-level
	representations.
	
	\section{Attention Entropy}
	
	For attention weights \(\{a_{ij}\}\),
	
	\[
	\Entropy_i^{\mathrm{attention}}
	=
	-\sum_j a_{ij}\log a_{ij}.
	\]
	
	\section{Patch Uncertainty}
	
	Patch uncertainty may combine prediction entropy, token similarity, and
	attention dispersion.
	
	\section{Semantic Entropy}
	
	Semantic entropy groups predictions according to semantic equivalence rather
	than exact output identity.
	
	%%%%%%%%%%%%%%%%%%%%%%%%%%%%%%%%%%%%%%%%%%%%%%%%%%%%%%%%%%%%%%%%%%%%%%%%%%%%%%%%
	% Chapter 16
	%%%%%%%%%%%%%%%%%%%%%%%%%%%%%%%%%%%%%%%%%%%%%%%%%%%%%%%%%%%%%%%%%%%%%%%%%%%%%%%%
	
	\chapter{Entropy in Foundation Models}
	\label{chap:foundation-entropy}
	
	\begin{chapteroverview}
		This chapter investigates uncertainty in large pretrained vision and
		vision-language models.
	\end{chapteroverview}
	
	\section{Segment Anything}
	
	Entropy can be computed over candidate masks, mask confidences, prompt
	responses, and dense image embeddings.
	
	\section{CLIP}
	
	Entropy may characterize image--text alignment, class-prompt uncertainty, and
	zero-shot prediction distributions.
	
	\section{DINO}
	
	Self-supervised transformer representations can be analyzed using token,
	feature, covariance, and neighborhood entropy.
	
	\section{SAM2}
	
	Video segmentation introduces temporal entropy, memory uncertainty, and
	mask-propagation uncertainty.
	
	\section{Vision-Language Models}
	
	Vision-language models require uncertainty measures that account for both
	visual representations and generated language.
	
	%%%%%%%%%%%%%%%%%%%%%%%%%%%%%%%%%%%%%%%%%%%%%%%%%%%%%%%%%%%%%%%%%%%%%%%%%%%%%%%%
	% Chapter 17
	%%%%%%%%%%%%%%%%%%%%%%%%%%%%%%%%%%%%%%%%%%%%%%%%%%%%%%%%%%%%%%%%%%%%%%%%%%%%%%%%
	
	\chapter{Entropy-Guided Learning}
	\label{chap:entropy-guided-learning}
	
	\begin{chapteroverview}
		This chapter explains how entropy can be incorporated directly into model
		training, sample selection, curriculum design, and explanation.
	\end{chapteroverview}
	
	\section{Entropy Loss}
	
	A general entropy-regularized objective may be written as
	
	\[
	\mathcal{L}
	=
	\mathcal{L}_{\mathrm{task}}
	+
	\lambda\mathcal{L}_{\mathrm{entropy}}.
	\]
	
	\section{Active Learning}
	
	Samples with high predictive or region-level entropy can be prioritized for
	annotation.
	
	\section{Curriculum Learning}
	
	Training may begin with low-entropy examples and gradually introduce more
	ambiguous examples.
	
	\section{Semi-Supervised Learning}
	
	Entropy minimization can encourage confident predictions on unlabeled data:
	
	\[
	\mathcal{L}_{\mathrm{ent}}
	=
	-\frac{1}{N}
	\sum_{i=1}^{N}
	\sum_{c=1}^{C}
	p_{ic}\log p_{ic}.
	\]
	
	\section{Explainable Artificial Intelligence}
	
	Entropy maps can indicate which regions, channels, tokens, or predictions are
	responsible for uncertainty.
	
	%%%%%%%%%%%%%%%%%%%%%%%%%%%%%%%%%%%%%%%%%%%%%%%%%%%%%%%%%%%%%%%%%%%%%%%%%%%%%%%%
	% PART V
	%%%%%%%%%%%%%%%%%%%%%%%%%%%%%%%%%%%%%%%%%%%%%%%%%%%%%%%%%%%%%%%%%%%%%%%%%%%%%%%%
	
	\part{Applications}
	
	%%%%%%%%%%%%%%%%%%%%%%%%%%%%%%%%%%%%%%%%%%%%%%%%%%%%%%%%%%%%%%%%%%%%%%%%%%%%%%%%
	% Chapter 18
	%%%%%%%%%%%%%%%%%%%%%%%%%%%%%%%%%%%%%%%%%%%%%%%%%%%%%%%%%%%%%%%%%%%%%%%%%%%%%%%%
	
	\chapter{Medical Imaging}
	\label{chap:medical-imaging}
	
	\begin{chapteroverview}
		Medical imaging presents uncertainty caused by noise, weak boundaries,
		biological variation, limited annotations, and disagreement among experts.
	\end{chapteroverview}
	
	\section{Brain MRI}
	
	Entropy methods may support tumor segmentation, tissue classification,
	registration, and uncertainty estimation.
	
	\section{Computed Tomography}
	
	Applications include liver segmentation, lesion detection, organ delineation,
	and low-dose reconstruction.
	
	\section{Histopathology}
	
	Entropy may describe cellular texture, tissue heterogeneity, and uncertainty in
	patch-level classification.
	
	\section{Retinal Imaging}
	
	Applications include retinal vessel segmentation and lesion detection.
	
	\section{Skin Lesions}
	
	Fuzzy, rough, probabilistic, and deep entropy may help quantify uncertain lesion
	boundaries.
	
	\section*{Representative Datasets}
	\addcontentsline{toc}{section}{Representative Datasets}
	
	\begin{datasetlist}
		\item BraTS;
		\item ISIC;
		\item DRIVE;
		\item LiTS.
	\end{datasetlist}
	
	%%%%%%%%%%%%%%%%%%%%%%%%%%%%%%%%%%%%%%%%%%%%%%%%%%%%%%%%%%%%%%%%%%%%%%%%%%%%%%%%
	% Chapter 19
	%%%%%%%%%%%%%%%%%%%%%%%%%%%%%%%%%%%%%%%%%%%%%%%%%%%%%%%%%%%%%%%%%%%%%%%%%%%%%%%%
	
	\chapter{Remote Sensing}
	\label{chap:remote-sensing}
	
	\begin{chapteroverview}
		Remote sensing combines spatial, spectral, temporal, and sensor uncertainty.
	\end{chapteroverview}
	
	\section{Hyperspectral Imaging}
	
	Hyperspectral images contain hundreds of correlated spectral bands and often
	require entropy-based feature selection.
	
	\section{Synthetic Aperture Radar}
	
	SAR imaging introduces speckle noise and uncertainty associated with scattering
	mechanisms.
	
	\section{Satellite Segmentation}
	
	Entropy can support land-cover segmentation, change detection, and active
	learning.
	
	\section*{Representative Datasets}
	\addcontentsline{toc}{section}{Representative Datasets}
	
	\begin{datasetlist}
		\item Indian Pines;
		\item Pavia University;
		\item Salinas.
	\end{datasetlist}
	
	%%%%%%%%%%%%%%%%%%%%%%%%%%%%%%%%%%%%%%%%%%%%%%%%%%%%%%%%%%%%%%%%%%%%%%%%%%%%%%%%
	% Chapter 20
	%%%%%%%%%%%%%%%%%%%%%%%%%%%%%%%%%%%%%%%%%%%%%%%%%%%%%%%%%%%%%%%%%%%%%%%%%%%%%%%%
	
	\chapter{Autonomous Driving}
	\label{chap:autonomous-driving}
	
	\begin{chapteroverview}
		Autonomous driving requires uncertainty-aware perception under changing
		illumination, weather, traffic, and sensor conditions.
	\end{chapteroverview}
	
	\section{Road Segmentation}
	
	Entropy maps identify uncertain road boundaries and difficult environmental
	conditions.
	
	\section{Lane Detection}
	
	Lane uncertainty arises from occlusion, worn markings, shadows, and weather.
	
	\section{Object Segmentation}
	
	Instance- and semantic-segmentation uncertainty can be measured at pixel,
	object, and scene levels.
	
	\section{Scene Understanding}
	
	Graph and region entropy may represent uncertainty in interactions among roads,
	vehicles, pedestrians, and infrastructure.
	
	\section*{Representative Datasets}
	\addcontentsline{toc}{section}{Representative Datasets}
	
	\begin{datasetlist}
		\item Cityscapes;
		\item BDD100K.
	\end{datasetlist}
	
	%%%%%%%%%%%%%%%%%%%%%%%%%%%%%%%%%%%%%%%%%%%%%%%%%%%%%%%%%%%%%%%%%%%%%%%%%%%%%%%%
	% Chapter 21
	%%%%%%%%%%%%%%%%%%%%%%%%%%%%%%%%%%%%%%%%%%%%%%%%%%%%%%%%%%%%%%%%%%%%%%%%%%%%%%%%
	
	\chapter{Industrial Vision}
	\label{chap:industrial-vision}
	
	\begin{chapteroverview}
		Industrial vision uses entropy for defect detection, quality assessment, and
		uncertainty-aware inspection.
	\end{chapteroverview}
	
	\section{Surface Inspection}
	
	Local texture entropy can reveal scratches, cracks, contamination, and
	irregular patterns.
	
	\section{Defect Detection}
	
	Feature and latent entropy can help distinguish rare defects from normal
	variation.
	
	\section{Quality Control}
	
	Entropy-based confidence estimates can determine when an automatic system
	should defer to human inspection.
	
	\section*{Representative Datasets}
	\addcontentsline{toc}{section}{Representative Datasets}
	
	\begin{datasetlist}
		\item MVTec AD.
	\end{datasetlist}
	
	%%%%%%%%%%%%%%%%%%%%%%%%%%%%%%%%%%%%%%%%%%%%%%%%%%%%%%%%%%%%%%%%%%%%%%%%%%%%%%%%
	% PART VI
	%%%%%%%%%%%%%%%%%%%%%%%%%%%%%%%%%%%%%%%%%%%%%%%%%%%%%%%%%%%%%%%%%%%%%%%%%%%%%%%%
	
	\part{Unified Framework}
	
	%%%%%%%%%%%%%%%%%%%%%%%%%%%%%%%%%%%%%%%%%%%%%%%%%%%%%%%%%%%%%%%%%%%%%%%%%%%%%%%%
	% Chapter 22
	%%%%%%%%%%%%%%%%%%%%%%%%%%%%%%%%%%%%%%%%%%%%%%%%%%%%%%%%%%%%%%%%%%%%%%%%%%%%%%%%
	
	\chapter{A Unified Entropy Framework}
	\label{chap:unified-framework}
	
	\begin{chapteroverview}
		This chapter unifies representations, uncertainty models, entropy measures,
		learning architectures, and applications.
	\end{chapteroverview}
	
	\section{Evolution of Representations}
	
	\begin{center}
		\begin{tikzpicture}[
			node distance=7mm,
			stage/.style={
				draw,
				rounded corners,
				minimum width=6.8cm,
				minimum height=8mm,
				align=center,
				fill=green!6
			},
			arrow/.style={
				-{Latex[length=2.5mm]},
				thick
			}
			]
			
			\node[stage] (p1) {Pixels};
			\node[stage, below=of p1] (p2) {Local neighborhoods and textures};
			\node[stage, below=of p2] (p3) {Regions and superpixels};
			\node[stage, below=of p3] (p4) {Graphs and structural representations};
			\node[stage, below=of p4] (p5) {Deep features and latent embeddings};
			\node[stage, below=of p5] (p6) {Foundation-model representations};
			
			\draw[arrow] (p1) -- (p2);
			\draw[arrow] (p2) -- (p3);
			\draw[arrow] (p3) -- (p4);
			\draw[arrow] (p4) -- (p5);
			\draw[arrow] (p5) -- (p6);
			
		\end{tikzpicture}
	\end{center}
	
	\section{Evolution of Uncertainty Models}
	
	\begin{center}
		\begin{tikzpicture}[
			node distance=7mm,
			stage/.style={
				draw,
				rounded corners,
				minimum width=6.8cm,
				minimum height=8mm,
				align=center,
				fill=orange!8
			},
			arrow/.style={
				-{Latex[length=2.5mm]},
				thick
			}
			]
			
			\node[stage] (u1) {Probability};
			\node[stage, below=of u1] (u2) {Fuzzy sets};
			\node[stage, below=of u2] (u3) {Rough sets};
			\node[stage, below=of u3] (u4) {Fuzzy-rough sets};
			\node[stage, below=of u4] (u5) {Evidence theory};
			\node[stage, below=of u5] (u6) {Learned uncertainty};
			
			\draw[arrow] (u1) -- (u2);
			\draw[arrow] (u2) -- (u3);
			\draw[arrow] (u3) -- (u4);
			\draw[arrow] (u4) -- (u5);
			\draw[arrow] (u5) -- (u6);
			
		\end{tikzpicture}
	\end{center}
	
	\section{Evolution of Entropy Measures}
	
	\begin{table}[htbp]
		\centering
		\caption{Entropy measures and their underlying representations.}
		\label{tab:entropy-representations}
		
		\begin{tabularx}{\textwidth}{
				>{\raggedright\arraybackslash}p{0.29\textwidth}
				>{\raggedright\arraybackslash}X
			}
			\toprule
			Representation &
			Representative entropy measures
			\\
			\midrule
			
			Pixel histograms &
			Shannon, Rényi, Tsallis
			\\
			
			Local neighborhoods &
			Local entropy, approximate entropy, sample entropy
			\\
			
			Textures &
			GLCM entropy, wavelet entropy, multiresolution entropy
			\\
			
			Regions &
			Region entropy, partition entropy
			\\
			
			Graphs &
			Structural entropy, spectral entropy, random-walk entropy
			\\
			
			Fuzzy memberships &
			Fuzzy entropy, intuitionistic fuzzy entropy
			\\
			
			Rough approximations &
			Rough entropy, neighborhood rough entropy
			\\
			
			Fuzzy-rough structures &
			Fuzzy-rough entropy
			\\
			
			Belief masses &
			Evidence entropy, Deng entropy
			\\
			
			Deep representations &
			Feature entropy, latent entropy, attention entropy
			\\
			
			\bottomrule
		\end{tabularx}
	\end{table}
	
	\section{Evolution of Learning Models}
	
	The relevant learning paradigms include:
	
	\begin{enumerate}
		\item classical image-processing algorithms;
		\item statistical machine learning;
		\item convolutional neural networks;
		\item graph neural networks;
		\item Vision Transformers;
		\item self-supervised models;
		\item vision foundation models;
		\item multimodal foundation models.
	\end{enumerate}
	
	\section{Unified Mathematical Formulation}
	
	Let
	
	\[
	\mathcal{R}(x;\theta)
	\]
	
	denote a representation of input \(x\), where \(\theta\) may contain fixed
	parameters, learned parameters, or both.
	
	Let
	
	\[
	\mathcal{U}
	\bigl(
	\mathcal{R}(x;\theta)
	\bigr)
	\]
	
	denote an uncertainty structure constructed on that representation.
	
	A general entropy functional may then be written as
	
	\[
	\boxed{
		\Entropy
		=
		\Entropy
		\left[
		\mathcal{U}
		\left(
		\mathcal{R}(x;\theta)
		\right)
		\right].
	}
	\]
	
	The representation may consist of:
	
	\begin{itemize}
		\item pixel values;
		\item local neighborhoods;
		\item image regions;
		\item graph nodes and edges;
		\item CNN feature maps;
		\item transformer tokens;
		\item latent embeddings.
	\end{itemize}
	
	The uncertainty model may be:
	
	\begin{itemize}
		\item probabilistic;
		\item fuzzy;
		\item rough;
		\item fuzzy-rough;
		\item evidence-theoretic;
		\item learned.
	\end{itemize}
	
	\section{Local and Global Entropy Optimization}
	
	For segmentation, local entropy and global representation entropy may have
	different objectives.
	
	Local neighborhood uncertainty should often be minimized:
	
	\[
	\min_{\theta}
	\Entropy_{\mathrm{local}}.
	\]
	
	At the same time, global representation diversity or discriminability may be
	maximized:
	
	\[
	\max_{\theta}
	\Entropy_{\mathrm{global}}.
	\]
	
	A combined objective is
	
	\[
	\mathcal{L}
	=
	\mathcal{L}_{\mathrm{seg}}
	+
	\lambda_{\mathrm{local}}
	\Entropy_{\mathrm{local}}
	-
	\lambda_{\mathrm{global}}
	\Entropy_{\mathrm{global}}.
	\]
	
	A multiscale extension is
	
	\[
	\mathcal{L}
	=
	\mathcal{L}_{\mathrm{seg}}
	+
	\lambda_p\Entropy_{\mathrm{pixel}}
	+
	\lambda_r\Entropy_{\mathrm{region}}
	+
	\lambda_g\Entropy_{\mathrm{graph}}
	-
	\lambda_d\Entropy_{\mathrm{dataset}}.
	\]
	
	Here:
	
	\begin{itemize}
		\item pixel entropy promotes locally confident predictions;
		\item region entropy promotes coherent segmented regions;
		\item graph entropy promotes structural consistency;
		\item dataset-level entropy prevents representational collapse and
		preserves class diversity.
	\end{itemize}
	
	%%%%%%%%%%%%%%%%%%%%%%%%%%%%%%%%%%%%%%%%%%%%%%%%%%%%%%%%%%%%%%%%%%%%%%%%%%%%%%%%
	% Chapter 23
	%%%%%%%%%%%%%%%%%%%%%%%%%%%%%%%%%%%%%%%%%%%%%%%%%%%%%%%%%%%%%%%%%%%%%%%%%%%%%%%%
	
	\chapter{Future Directions}
	\label{chap:future-directions}
	
	\begin{chapteroverview}
		This chapter identifies emerging combinations of classical entropy theory,
		rough and fuzzy models, graphs, deep learning, and foundation models.
	\end{chapteroverview}
	
	\section{Deep Rough Entropy}
	
	Deep rough entropy applies rough approximations to CNN or transformer
	embeddings rather than raw attributes.
	
	A possible pipeline is:
	
	\begin{center}
		\begin{tikzpicture}[
			node distance=7mm,
			box/.style={
				draw,
				rounded corners,
				minimum width=7cm,
				minimum height=8mm,
				align=center,
				fill=blue!5
			},
			arrow/.style={
				-{Latex[length=2.5mm]},
				thick
			}
			]
			
			\node[box] (image) {Input image};
			\node[box, below=of image] (encoder)
			{CNN or Vision Transformer encoder};
			\node[box, below=of encoder] (embedding)
			{Pixel, patch, or region embeddings};
			\node[box, below=of embedding] (neighborhood)
			{Neighborhood or similarity relation};
			\node[box, below=of neighborhood] (rough)
			{Lower, upper, and boundary approximations};
			\node[box, below=of rough] (entropy)
			{Deep rough entropy};
			\node[box, below=of entropy] (objective)
			{Regularization, refinement, or uncertainty estimation};
			
			\draw[arrow] (image) -- (encoder);
			\draw[arrow] (encoder) -- (embedding);
			\draw[arrow] (embedding) -- (neighborhood);
			\draw[arrow] (neighborhood) -- (rough);
			\draw[arrow] (rough) -- (entropy);
			\draw[arrow] (entropy) -- (objective);
			
		\end{tikzpicture}
	\end{center}
	
	\section{Graph Rough Entropy}
	
	Graph rough entropy defines approximations over graph neighborhoods,
	communities, or diffusion structures.
	
	\section{Fuzzy-Rough Deep Learning}
	
	Fuzzy-rough deep learning combines:
	
	\begin{itemize}
		\item fuzzy similarity between learned embeddings;
		\item fuzzy lower and upper approximations;
		\item differentiable fuzzy-rough entropy;
		\item end-to-end neural optimization.
	\end{itemize}
	
	\section{Multiscale Entropy}
	
	A multiscale system may jointly measure uncertainty at:
	
	\begin{enumerate}
		\item pixel level;
		\item local-neighborhood level;
		\item region level;
		\item graph level;
		\item embedding level;
		\item dataset level.
	\end{enumerate}
	
	\section{Entropy-Aware Foundation Models}
	
	Research questions include:
	
	\begin{itemize}
		\item uncertainty of prompt-conditioned masks;
		\item entropy of dense foundation-model embeddings;
		\item token- and region-level calibration;
		\item cross-modal semantic uncertainty;
		\item reliability under domain shift.
	\end{itemize}
	
	\section{Uncertainty-Aware Artificial Intelligence}
	
	Entropy may support:
	
	\begin{itemize}
		\item selective prediction;
		\item human-in-the-loop decision-making;
		\item active learning;
		\item model calibration;
		\item out-of-distribution detection;
		\item explainability;
		\item safety-critical vision.
	\end{itemize}
	
	\section{Open Research Challenges}
	
	\begin{researchquestions}
		\item How should entropy be defined for high-dimensional deep embeddings?
		\item How can rough and fuzzy approximations be made fully differentiable?
		\item Which entropy objectives improve segmentation without causing
		over-smoothing?
		\item How should local entropy minimization be balanced against global
		entropy maximization?
		\item How can entropy measures be evaluated independently of downstream
		accuracy?
		\item How should uncertainty be calibrated across pixels, regions,
		graphs, and complete images?
		\item Can one unified entropy measure operate consistently across
		multiple representation levels?
	\end{researchquestions}
	
	%%%%%%%%%%%%%%%%%%%%%%%%%%%%%%%%%%%%%%%%%%%%%%%%%%%%%%%%%%%%%%%%%%%%%%%%%%%%%%%%
	% Chapter 24
	%%%%%%%%%%%%%%%%%%%%%%%%%%%%%%%%%%%%%%%%%%%%%%%%%%%%%%%%%%%%%%%%%%%%%%%%%%%%%%%%
	
	\chapter{Conclusions}
	\label{chap:conclusions}
	
	Entropy in image processing has evolved from a statistic computed over
	intensity histograms into a broad family of measures operating on local
	neighborhoods, textures, regions, graphs, feature maps, attention patterns,
	and latent representations.
	
	At the same time, models of uncertainty have expanded from probability to
	fuzzy sets, rough sets, fuzzy-rough approximations, evidence theory, Bayesian
	deep learning, and learned uncertainty.
	
	The unified perspective developed in this book can be summarized as
	
	\[
	\boxed{
		\text{representation}
		+
		\text{uncertainty model}
		+
		\text{entropy functional}
		+
		\text{learning objective}
		+
		\text{application}.
	}
	\]
	
	This viewpoint supports both systematic comparison of existing methods and the
	development of new methods such as deep rough entropy, graph rough entropy,
	fuzzy-rough representation learning, multiscale entropy optimization, and
	entropy-aware foundation models.
	
	%%%%%%%%%%%%%%%%%%%%%%%%%%%%%%%%%%%%%%%%%%%%%%%%%%%%%%%%%%%%%%%%%%%%%%%%%%%%%%%%
	% Appendices
	%%%%%%%%%%%%%%%%%%%%%%%%%%%%%%%%%%%%%%%%%%%%%%%%%%%%%%%%%%%%%%%%%%%%%%%%%%%%%%%%
	
	\appendix
	
	\chapter{Mathematical Background}
	\label{app:mathematical-background}
	
	\section{Probability Distributions}
	
	A discrete probability distribution satisfies
	
	\[
	p_i\geq0,
	\qquad
	\sum_{i=1}^{n}p_i=1.
	\]
	
	\section{Similarity Relations}
	
	A similarity relation may be defined using a distance \(d\):
	
	\[
	R_{\sigma}(x,y)
	=
	\exp
	\left(
	-\frac{d(x,y)^2}{2\sigma^2}
	\right).
	\]
	
	\section{Graph Construction}
	
	Given embeddings \(\bm{z}_1,\ldots,\bm{z}_n\), a weighted graph may use
	
	\[
	w_{ij}
	=
	\exp
	\left(
	-\frac{
		\lVert\bm{z}_i-\bm{z}_j\rVert_2^2
	}{
		2\sigma^2
	}
	\right).
	\]
	
	\chapter{Experimental Protocol Templates}
	\label{app:experimental-protocols}
	
	\section{Segmentation Protocol}
	
	A standard segmentation experiment should specify:
	
	\begin{itemize}
		\item dataset;
		\item train, validation, and test splits;
		\item architecture;
		\item preprocessing;
		\item augmentation;
		\item optimizer;
		\item learning-rate schedule;
		\item entropy formulation;
		\item regularization coefficient;
		\item evaluation metrics;
		\item statistical testing.
	\end{itemize}
	
	\section{Segmentation Metrics}
	
	Intersection over Union is
	
	\[
	\operatorname{IoU}
	=
	\frac{
		|\widehat{Y}\cap Y|
	}{
		|\widehat{Y}\cup Y|
	}.
	\]
	
	The Dice coefficient is
	
	\[
	\operatorname{Dice}
	=
	\frac{
		2|\widehat{Y}\cap Y|
	}{
		|\widehat{Y}|+|Y|
	}.
	\]
	
	\section{Uncertainty Metrics}
	
	Recommended uncertainty metrics include:
	
	\begin{itemize}
		\item expected calibration error;
		\item negative log-likelihood;
		\item Brier score;
		\item area under the risk--coverage curve;
		\item boundary uncertainty accuracy;
		\item out-of-distribution detection performance.
	\end{itemize}
	
	%%%%%%%%%%%%%%%%%%%%%%%%%%%%%%%%%%%%%%%%%%%%%%%%%%%%%%%%%%%%%%%%%%%%%%%%%%%%%%%%
	% Back matter
	%%%%%%%%%%%%%%%%%%%%%%%%%%%%%%%%%%%%%%%%%%%%%%%%%%%%%%%%%%%%%%%%%%%%%%%%%%%%%%%%
	
	\backmatter
	
	\chapter*{Bibliography}
	\addcontentsline{toc}{chapter}{Bibliography}
	
	A BibTeX bibliography can be enabled by replacing this placeholder with:
	
	\begin{quote}
		\texttt{\textbackslash bibliographystyle\{plain\}}\\
		\texttt{\textbackslash bibliography\{references\}}
	\end{quote}
	
	Alternatively, the document may use the \texttt{biblatex} package and
	\texttt{\textbackslash printbibliography}.
	
	\chapter*{Index}
	\addcontentsline{toc}{chapter}{Index}
	
	An index may be generated later using the \texttt{makeidx} or
	\texttt{imakeidx} package.
	
\end{document}